\documentclass{article}
\usepackage[UTF8]{ctex}
\usepackage{amsmath,amssymb,amsfonts}
\usepackage{enumitem}
\usepackage{geometry}
\geometry{a4paper, margin=1in}
% === 错题本排版优化（不参与题目标识） ===
\setlength{\parindent}{0pt}
\setlength{\parskip}{0.35em plus 0.10em minus 0.08em}
\setlength{\emergencystretch}{3em}
\setlist{itemsep=0.25em, topsep=0.40em, parsep=0pt}
\allowdisplaybreaks[2]
% === 错题本排版优化结束 ===
\begin{document}
\section*{中国人民大学2025年数学分析真题}
\section*{20．中国人民大学2025年数学分析试卷}

\subsection{题目1}
1．求最小的实数 \(a\) ，使得对任意的 \(x>0\) ，都有 \(\left(1+\frac{1}{x}\right)^{x+a}>e\) 。

\subsection{题目2}
2．已知 \(x_{0}, y_{0}, z_{0}\) 为三个给定的实数，对正整数 \(n \geq 1\) ，令

\[
x_{n}=\frac{y_{n-1}+z_{n-1}}{2}, y_{n}=\frac{x_{n-1}+z_{n-1}}{2}, z_{n}=\frac{x_{n-1}+y_{n-1}}{2} .
\]

证明：序列 \(\left\{x_{n}\right\},\left\{y_{n}\right\},\left\{z_{n}\right\}\) 均收敛，并计算它们的极限．

\subsection{题目3}
3．设函数 \(f(x)\) 在 \([0,1]\) 上连续，在 \((0,1)\) 上可导，且 \(f(0)=0, f(1)=2\) ．证明：存在 \(\xi, \eta \in(0,1), \xi \neq\) \(\eta\) ，满足 \(\left[1+f^{\prime}(\xi)\right]\left[1+f^{\prime}(\eta)\right]=9\) ．

\subsection{题目4}
4．设 \(n\) 为正整数，计算极限 \(\lim _{n \rightarrow \infty} \int_{0}^{1} \sqrt[n]{x^{n}+(1-x)^{n}} \mathrm{~d} x\) ．

\subsection{题目5}
5．求三元函数 \(f\left(x_{1}, x_{2}, x_{3}\right)=\sum_{i=1}^{3} \sum_{j=1}^{3}\left(e^{x_{i}}-x_{j}\right)^{2}\) 的最小值，其中 \(x_{1}, x_{2}, x_{3}\) 均为实数．

\subsection{题目6}
6．已知定义在 \([0,3]\) 上的非负连续函数 \(f(x)\) 满足 \(\int_{0}^{3} \frac{1}{1+f(x)} \mathrm{d} x=1\) ，证明： \(\int_{0}^{3} \frac{f(x)}{1+f^{2}(x)} \mathrm{d} x \leq 1\) ．

\subsection{题目7}
7．设 \(D\) 是两条直线 \(y=x, y=4 x\) 和两条双曲线 \(x y=1, x y=4\) 在第一象限围成的区域，边界 \(\partial D\) 的方向为逆时针方向，\(F(u)\) 为 \(\mathbb{R}\) 上的可导函数，\(F^{\prime}(u)=f(u)\) ，且 \(f(u)\) 在 \(\mathbb{R}\) 中连续，证明：

\[
\int_{\partial D} \frac{F(x y)}{y} \mathrm{~d} y=\ln 2 \int_{1}^{4} f(u) \mathrm{d} u .
\]

\subsection{题目8}
8．设 \(\sum\) 为上半椭球面 \(\frac{x^{2}}{2}+\frac{y^{2}}{2}+z^{2}=1, z \geq 0\) 。对 \(\sum\) 上任意一点 \((x, y, z)\) ，记 \(\rho(x, y, z)\) 为原点到 \(\sum\)在 \((x, y, z)\) 处的切平面的距离，计算 \(\iint_{\sum} \frac{z}{\rho(x, y, z)} \mathrm{d} S\) ．

\subsection{题目9}
9．已知对任意的正整数 \(n, f_{n}(x)\) 均为 \([a, b]\) 上的单调函数，且 \(\left\{f_{n}(x)\right\}\) 在 \([a, b]\) 上逐点收敛于连续函数 \(f(x)\) ．证明：\(\left\{f_{n}(x)\right\}\) 在 \([a, b]\) 上一致收敛于 \(f(x)\) ．

\subsection{题目10}
10．对正整数 \(n\) ，判断反常积分 \(\int_{0}^{+\infty} x \sin \left(x^{n}\right) \mathrm{d} x\) 是否收敛以及是否绝对收敛，并说明理由．
\end{document}
